\section{What the evidence establishes}
\label{sec:results}

\subsection{A result belongs to a cell and a revision}

\readcode The status-detail document warns that much of its main body is September history.
The README and newer per-topic records add Windows, display and relay results. This paper
therefore does not copy a historical heading such as ``Windows: research only'' into a current
support claim. Conversely, a later source audit does not rerun an earlier hardware test.
The following table deliberately retains the recorded revisions.\src{status}\src{readme}

\begin{center}\small
\begin{tabularx}{\textwidth}{@{}p{25mm}p{27mm}X@{}}
\toprule Cell & Revision & Recorded scope; repetition limit \\
\midrule
Linux / TU116 & \code{1915bd71} & Gates 9/9, bare and thin 30/30; one recorded merge-bar
campaign. Ladder at \code{a295e2ce}, described as the same code. No timing spread claimed.
\src{turing} \\
Linux / AD104 workstation & \code{fc2dadc4} & Gates 9/9, bare and thin 30/30; ladder
4/4 at \code{d40dafbf}. Distinct checkout revisions are preserved.\src{families} \\
Linux / GB205 & \code{d40dafbf} & Gates 9/9, bare and thin 30/30, ladder 4/4.
One recorded cell, not a Windows pass.\src{families} \\
Linux / GA104 & \code{714ec29d3} & Host and guest 595.91.07; gates 9/9 plus gate 10,
bare and thin 30/30, one box/cell.\src{requirements}\src{sweep} \\
Windows / AD104 desktop & \code{eb9f2aa3} & Stock 580.88 application matrix; later
relay session at \code{89fc8a222}. Completion, recovery and output grades differ by app.
No full-matrix pass at the paper baseline is asserted.\src{readme}\src{winmatrix} \\
\bottomrule
\end{tabularx}
\end{center}

\measured The family records above are dated 2026-09-30; the GA104 record is dated
2026-10-24. Pass counts enumerate tests or arms, not independent boots of the entire workload.
A 30/30 thin suite boots separate arms; it is not a thirty-boot confidence interval for a
multi-application desktop. The Linux family campaigns used nested rented hosts and open
580.159.04 drivers unless the cell says otherwise.\src{turing}\src{families}\src{sweep}

\readcode The public requirements list additional earlier Linux runs on GA106, GA102,
AD106 and GB203, and the README mentions an AD102 Linux guest with driver 590. Windows desktops
on AD106 and AD103 used a borrowed display page. The requirements explicitly say that no single
current per-die table is documented. This is evidence fragmentation, not permission to call
all neighbouring dies supported. Hopper and datacentre families remain source-derived or
refused cases, not hardware qualifications.\src{requirements}\src{readme}

\subsection{Correctness results and their sample units}

\measured The recorded Linux application matrix reports 61/65 passing applications and six
stream probes, with output comparisons to the same host; its source is an earlier GA106
campaign, not the paper baseline. The four remaining failures need managed-memory demand
paging. The record of 100 sequential processes applies to one persistence-mode boot. These
are useful workload and lifetime samples, but neither is a soak distribution.\src{status}
\src{apps}\src{readme}

\measured The headless graphics record reports 38/38 items, including 31 byte-identical
outputs, at \code{d06833f0} on a GA104 box. The video lane reports byte-identical encode/decode
outputs on a GA106. These compare content, which is stronger than a renderer name; they do
not cover every display mode, all video families or later revisions.\src{status}\src{video}

\readcode The driver matrix varies host and guest independently. Several guest versions
passed a CUDA ladder against a fixed host, and several hosts passed the gates, thin suite and
ladder against a fixed guest. The public guide preserves the versions and revisions and warns
that the raw 30-arm suite is graded for 580.x guests. A decoder's support range and a tested
application range are different sets.\src{requirements}\src{driverdesign}

\subsection{A performance row with actual spread}

\measured At device revision \code{e5ff45ff}, the AMD EPYC 7452 nested host with RTX 3060
and 580.159.04 ran eager Qwen2-0.5B decoding. For 512-token warm decode, the summary gives
\textbf{12.62 tokens/s [12.40, 12.94], $n=3$} in the guest and
\textbf{43.98 [43.59, 44.63], $n=3$} on the host, ratio \textbf{0.287}.
These are the table's reported aggregates and ranges, not a new estimate or confidence
interval. The associated text-equality checks matter separately from throughput.
\src{llmmethod}\src{llmdata}

\measured The Windows performance summary, 2026-10-17, records three experimental
Direct3D 11 Heaven runs at \code{134dcb54}: \textbf{299.1, 296.5 and 299.2 fps}
with scan-off and lazy observation together, 1920$\times$1080, Ultra, 260-second windows,
stock 580.88 guest on the owner's AD104 host with 595.91.07. Its production-path samples at
other listed revisions are slower. The record itself says these experiments were not a
safe integrated lifecycle qualification and did not have an application matrix. Their
combined gain cannot be assigned to just one optimisation.\src{winperf}

\inferred These examples support two narrower conclusions: large batches and small launches
need separate measurements, and a fast renderer can coexist with a correctness or latency
problem under mixed use. They do not establish a universal ``virtualisation overhead'' number.
Nested versus non-nested execution, host load, guest workload and baseline revision all change
the denominator.\src{linuxperf}\src{winperf}

\subsection{The Worker GPU gate: useful, unattested and owner-controlled}

\readcode GitHub requests a check through a Cloudflare Worker. A Durable Object ledger
reserves liability, tracks owned instances, deadlines and spending, and coordinates cleanup.
The provider account key stays in the Worker. Identity checks bind the workflow, owner
triggering actor, signed run attempt and tested SHA. A scoped box capability can report only
within its run contract. Cloudflare Access reduces noise; it is not the content-approval
boundary.\src{broker}\src{threat}

\readcode At the baseline's contributor policy, ordinary PRs target \code{next}; the GPU
check belongs on the exact final \code{next} to \code{master} batch. The policy states CI is
required on next and CI plus \code{gpu-vast} on master. This is the recorded policy, not an
audit of live GitHub ruleset settings. Older triage or handoff statements about required
checks are historical.\src{rules}

\attack{Only the owner-approved SHA rents within the configured policy, and cleanup acts on
ledger-owned instances.}
{Broker unit/runtime tests and recorded provider campaigns exercise identity, accounting,
results, retries and teardown. A real runtime matters because Node fakes cannot reproduce
all Worker fetch behaviour.\src{threat}\src{broker}}
{The rented machine and its operator can lie. A valid callback authenticates the run capability,
not honest GPU execution. A PASS is a regression tripwire, not release provenance, an
attestation or a security audit. A missed cron or failed provider deletion can also outlive
an internal deadline.\src{broker}\src{threat}}
